\section{A practical hierarchy for persistence and dependent inference}
\label{s5:response-inference}
A treasury analyst can reasonably want a deposit response at horizons zero
through twelve, a simultaneous statement about that whole curve, and an
answer robust to a weak monetary-policy instrument. These are three
requirements. Changing the covariance bandwidth addresses dependence;
changing to a reduced-form projection may change the target; neither action
establishes the instrument's exclusion restriction.

For the direct-IV route, stack the raw moments as
\[
 m_t=Z_t\big((X_{i,t+h})_{(i,h)\in\mathcal I\times\mathcal H}',X_{1t}\big)',
 \quad \mu=\E[m_t],\quad \theta_j=\mu_j/\mu_d,
 \quad r_t(\theta)=[I_{d-1},-\theta]m_t.
\]
Here $d-1$ is the number of response coordinates, $X_{1t}$ supplies the
impact normalization, and $j$ denotes an outcome-horizon pair. Relevance,
lead/lag exclusion and the structural linear response argument above make
these ratios causal. A fixed-parameter stationary CLT gives
$\sqrt n(\bar m-\mu)\Rightarrow N(0,\Omega)$.
For each coordinate, invert the residual-moment test
\begin{equation}
 n(\bar m_j-\vartheta\bar m_d)^2\le
 c^2\{\hat\Omega_{jj}-2\vartheta\hat\Omega_{jd}
                +\vartheta^2\hat\Omega_{dd}\}.
 \label{s5:residual-set}
\end{equation}
Using $c=\Phi^{-1}\{1-\alpha/[2(d-1)]\}$ gives a Bonferroni whole-curve
set when every marginal test has the asserted size. No division by
$\bar m_d$ is needed. The answer can be disconnected, unbounded, or empty;
a plotting system must retain that topology. This is a rectangular joint
construction, not an efficient ellipsoidal region or a uniform theorem over
all weakly dependent processes.

\subsection{What prefiltering assumes and what it changes}
The Bartlett long-run covariance estimator uses
$\hat\Gamma_0+\sum_{\ell=1}^L(1-\ell/(L+1))
(\hat\Gamma_\ell+\hat\Gamma_\ell')$ \citep{newey1986}. Its positive
semidefiniteness does not imply a good finite-sample approximation near a
unit root. For centered column scores $u_t=m_t-\mu$, define
$e_t=u_t-Au_{t-1}$. If the covariance sums exist absolutely,
\[
 \Omega_e=(I-A)\Omega_u(I-A)',\qquad
 \Omega_u=(I-A)^{-1}\Omega_e(I-A)^{-T}.
\]
To see this, sum $e_t$ over $t=1,\ldots,n$:
$(I-A)\sum_tu_t+A(u_n-u_0)$. The boundary term is $o_p(\sqrt n)$ in a
stationary finite-variance model. Taking limiting covariances gives the
identity. The residuals need not be white: a misspecified stable AR(1)
prefilter can still be followed by a consistent residual HAC estimator.

Feasible consistency needs more than this identity. For fixed dimension,
assume $\hat A\to A$, $I-A$ is uniformly nonsingular in the chosen fixed
model, and the HAC estimator using estimated residuals differs by $o_p(1)$
from the population-residual estimator. A sufficient replacement condition
under bounded sample quadratic moments is
$L\|\hat A-A\|=o_p(1)$; combine it with the HAC dependence/moment and
bandwidth conditions. Root-$n$ estimation of $A$ and $L=o(\sqrt n)$ give
one route. Continuity of matrix inversion then completes the argument.
An estimated radius below one is a useful numerical check, not evidence
of a uniform stability margin for the true bank process. Near singularity,
recolouring magnifies covariance error by the square of the inverse norm.

The implementation fits a VAR(1) to the score vector, estimates a residual
Bartlett covariance and uses linear solves for recolouring. It reports the
prefilter radius and condition number and refuses an unstable estimate.
The three compared methods use the \emph{same} origins, instruments,
normalization and test inversion. Only covariance construction changes.
This preserves the intervention while revealing the inferential sensitivity.

\begin{center}\small
\begin{tabular}{lrrr}\toprule
External-IV regime & $L=\lfloor n^{1/4}\rfloor$ & $4L$ & AR prefilter\\\midrule
Stationary & .940 & .928 & .935\\
Persistent instrument noise & .905 & .935 & .923\\
Near-unit dynamics & .758 & .933 & .775\\
Near-unit, longer series & .803 & .973 & .820\\
Weak instrument & .910 & .963 & .983\\
Invalid instrument & .240 & .363 & .273\\\bottomrule
\end{tabular}\end{center}
These are 400-repetition whole-curve coverages over 26 coordinates. Weak-set
coverage may improve by becoming unbounded, so coverage must be read with
the archived unbounded-set frequencies and failure diagnostics. In the weak
case, both longer-window and prefiltered constructions are unbounded in at
least one coordinate in 99.25\% of repetitions. The earlier
delta-method curve's 81.8\% near-unit coverage remains an adverse result;
it is not overwritten by this different residual-test comparison. The
longer window is a useful primary challenger for this declared stationary
external-IV family. Prefiltering is not uniformly better. No covariance
method repairs invalid exclusion, as the last row shows.

\subsection{Lag augmentation and a distinct joint construction}
For an AR(1), regress $y_{t+h}$ on $(1,y_t,y_{t-1})$. Partialling out
$y_{t-1}$ makes the effective regressor the innovation in $y_t$. Under the
source's strong innovation restrictions, the relevant score covariance
permits heteroskedasticity-robust inference without a generic residual HAC
penalty. \citet[Definition 1, Assumptions 2--3 and Proposition 1]{lpaug2020}
establishes a marginal result over a specified persistence/root class and
horizon-growth regime. For four prespecified horizons, applying that result
at level $\alpha/4$ and the union bound supplies a simultaneous route;
it does not require independent horizon estimates.

The scalar experiment uses horizons $(1,4,8,12)$, the Eicker--Huber--White
regression covariance and that Bonferroni critical value. Joint 95\%
coverage is .920, .943 and .928 for $\rho=.8,.985,1$ respectively, across
400 repetitions. This is evidence for an AR reduced-form response under
its innovation law. It is not a comparison on the external-IV normalized
target in Equation~\eqref{s5:residual-set}. Uniform marginal validity also
does not by itself permit a growing number of horizons without controlling
the accumulated error \citep{inoue2023}.

The instrument need not affect every date equally. In the inspected 2025
dynamic-LATE preprint, conditional instrument independence, exclusion,
monotonicity, partial first-stage support and smooth response restrictions
make a generalized Wald ratio a weighted marginal effect for the responsive
subpopulation \citep[Assumptions 2.1--2.5 and Proposition 2.1]{casini2025}.
Its further claim of date-specific complier recovery needs deterministic
complier status and local smoothness/comparison assumptions
\citep[Assumptions 2.6--2.8]{casini2025}. Neither statement licenses an
analyst to choose the dates with the strongest observed first stage and use
an ordinary full-sample interval. Selected subsamples change population
weights and require a separately justified selection-aware inferential argument.

\subsection{Choosing a route for a bank's actual clock and panel}
Common shocks require aggregation of customer scores by calendar date;
customer count does not manufacture independent shock realizations. In a
separate AR(.95) date-score mean experiment, short-window HAC covers .500
at 800 dates and .560 at 2,000 dates; prefiltering covers .918 and .958.
This favorable result supports prefiltering for that score model, while
the curve comparison above prevents generalizing it to every response.
With heavy tails, the moment and covariance conditions need checking or a
different robust sampling argument; trimming a shock changes the target
unless a retained-population estimand is specified.

The wider method hierarchy follows restrictions, not a package ranking.
Use direct external-IV projections when a dated external source of variation
and its exclusion are defensible. A finite VAR or factor reduction buys
precision and long-horizon propagation by imposing dynamics and rank; compare
it with the direct route on identical shock normalization. Use lag-augmented
projections for their justified reduced-form or separately identified
structural target. An error-correction model is appropriate when levels have
economically justified cointegrating restrictions; estimating rank and
deterministic trends changes the inferential problem and cannot inherit a
stationary score theorem. Differences answer change effects, with level
responses reconstructed by cumulation under the maintained integration
model. Mixed-frequency state-space/MIDAS designs add a measurement and
aggregation law: quarterly output is not three identical monthly observations.
The observation-clock and factor derivations earlier in this chapter explain
the added information and aliasing cost. Breaks require a declared stable
subperiod, a prespecified regime model or explicitly local effects; a
full-sample mean response can otherwise average incompatible policies.

These alternatives are proposed supported model families, not automatically
selected truths. The reusable response operation reports the stability,
normalization, dependence and aggregation premises, the primary estimate,
a substantive challenger and their disagreement. When no family is credible,
it returns an exploratory curve or a sensitivity set with that limitation.
Large-scale reduction remains valuable, but it follows the scientific unit
and target instead of choosing them.
